\section{Modeling Approach}



\subsection{Objective Function}
The objective is to maximize the total annual profit, which is the sum of revenues from energy arbitrage and ancillary service capacity payments, minus the cost of charging the battery.

\begin{align}
\text{Maximize} \quad & \sum_{t \in T} \left[ P_{DA}(t) \cdot p_{dis}(t) - P_{DA}(t) \cdot p_{ch}(t) \right] \cdot \Delta t \\
&+ \sum_{b \in B} \left[ P_{FCR}(b) \cdot c_{fcr}(b) + P_{aFRR}^{pos}(b) \cdot c_{afrr}^{pos}(b) + P_{aFRR}^{neg}(b) \cdot c_{afrr}^{neg}(b) \right] \\
&- \sum_{t \in T} \left[ p_{ch}(t) \cdot \frac{P_{DA}(t)}{1000} \right] \cdot \Delta t
\end{align}

\subsection{Constraints}
The model is subject to a series of constraints that define the physical and operational limits of the BESS:

\begin{enumerate}
    \item \textbf{State of Charge (SOC) Dynamics:} The energy level at the end of each interval is determined by the previous interval's energy level and the charge/discharge actions taken, adjusted for efficiency. For all $t \in T$:
    \begin{equation}
    e_{soc}(t) = e_{soc}(t-1) + \left[ p_{ch}(t) \cdot \eta_{ch} - \frac{p_{dis}(t)}{\eta_{dis}} \right] \cdot \Delta t
    \end{equation}

    \item \textbf{SOC Limits:} The stored energy must remain within the BESS's physical capacity. For all $t \in T$:
    \begin{equation}
    SOC_{min} \cdot E_{nom} \leq e_{soc}(t) \leq SOC_{max} \cdot E_{nom}
    \end{equation}

    \item \textbf{Co-optimization Power Limits:} These are the most critical constraints, linking the markets together. The power used for arbitrage plus the power reserved for ancillary services cannot exceed the BESS's configured maximum power. For each block $b \in B$ and each interval $t \in b$:
    \begin{align}
    p_{ch}(t) + c_{fcr}(b) \cdot 1000 + c_{afrr}^{pos}(b) \cdot 1000 &\leq P_{max}^{config} \\
    p_{dis}(t) + c_{fcr}(b) \cdot 1000 + c_{afrr}^{neg}(b) \cdot 1000 &\leq P_{max}^{config}
    \end{align}
    %% Question on the formulation above: why is it plus here? Shouldn't it be minus, as the capacity reserved for FCR and aFRR reduces the power available for arbitrage? <-- Ancillary service commitments in the opposite direction limit your arbitrage capacity in the current direction.
    These constraints force the model to make an economic trade-off: reserving 1 MW of capacity for FCR reduces the power available for energy arbitrage by 1 MW, potentially sacrificing arbitrage revenue for the certain income from the FCR capacity payment.

    \item \textbf{Simultaneous Operation Prevention:} The BESS cannot charge and discharge at the same time. This is enforced using binary variables and the "big-M" method. For all $t \in T$:
    \begin{align}
    p_{ch}(t) &\leq y_{ch}(t) \cdot P_{max}^{config} \\
    p_{dis}(t) &\leq y_{dis}(t) \cdot P_{max}^{config} \\
    y_{ch}(t) + y_{dis}(t) &\leq 1
    \end{align}

    \item \textbf{Daily Cycle Limit:} The total energy discharged over any 24-hour period (96 intervals) must not exceed the configured limit. For each day $d$ in the year:
    \begin{equation}
    \sum_{t \in \text{Day } d} p_{dis}(t) \cdot \Delta t \leq N_{cycles} \cdot E_{nom}
    \end{equation}

    \item \textbf{Ancillary Service Energy Reserve:} The BESS must maintain sufficient energy reserves to be able to deliver its awarded ancillary capacity for a specified duration (e.g., 15 minutes for FCR). For each block $b \in B$ and each interval $t \in b$:
    \begin{align}
    e_{soc}(t) &\geq SOC_{min} \cdot E_{nom} + \left( c_{fcr}(b) \cdot 1000 + c_{afrr}^{pos}(b) \cdot 1000 \right) \cdot 0.25 \\
    e_{soc}(t) &\leq SOC_{max} \cdot E_{nom} - \left( c_{fcr}(b) \cdot 1000 + c_{afrr}^{neg}(b) \cdot 1000 \right) \cdot 0.25
    \end{align}
\end{enumerate}

\subsection{Formulating the Co-Optimization Problem (MILP)}

The mathematical model aims to maximize total revenue by making optimal decisions for all variables across the entire time horizon simultaneously, subject to the physical and market constraints of the system.

\subsubsection{Sets and Indices:}
\begin{itemize}
    \item $t \in T$: The set of discrete 15-minute time intervals over the one-year simulation horizon. The duration of each interval is $\Delta t = 0.25$ hours.
    \item $b \in B$: The set of discrete 4-hour blocks for which ancillary service capacity bids are placed. Each block $b$ contains 16 time intervals $t$.
\end{itemize}

\subsubsection{Parameters:}
\begin{itemize}
    \item $P_{DA}(t)$: Day-Ahead electricity price for interval $t$.
    \item $P_{FCR}(b)$: FCR capacity price for block $b$.
    \item $P_{aFRR}^{pos}(b)$: Positive aFRR capacity price for block $b$.
    \item $P_{aFRR}^{neg}(b)$: Negative aFRR capacity price for block $b$.
    \item $E_{nom}$: Nominal energy capacity of the BESS (4,472 kWh).
    \item $P_{max}^{config}$: Maximum charge/discharge power based on the selected C-rate (0.25, 0.33, 0.50 C).
    \item $\eta_{ch}, \eta_{dis}$: BESS charging and discharging efficiencies (assumed constant, e.g., 95\% each (94.87\% exactly) for a round-trip efficiency of ~90\%).
    \item $SOC_{min}, SOC_{max}$: Minimum and maximum state of charge limits (0\% to 100\% of $E_{nom}$).
    \item $N_{cycles}$: The daily cycle limit for the selected configuration [FCE/day] (1, 1.5, 2 per day).
    % define the Initial state of the SoC
    \item $e_{soc}(0)$: Initial state of charge at the start of the simulation, $e_{soc}(0) \in [0, E_{nom}] $ (Notice that one do not need to ensure the final SOC the same as the starting SOC).
\end{itemize}

\subsubsection{Decision Variables:}
\begin{itemize}
    \item $p_{ch}(t) \geq 0$: Power used to charge the BESS from the grid during interval $t$.
    \item $p_{dis}(t) \geq 0$: Power discharged from the BESS to the grid during interval $t$.
    \item $e_{soc}(t) \geq 0$: Energy stored in the BESS at the end of interval $t$.
    \item $c_{fcr}(b) \geq 0$: Symmetric FCR capacity bid for block $b$.
    \item $c_{afrr}^{pos}(b) \geq 0$: Positive aFRR capacity bid for block $b$.
    \item $c_{afrr}^{neg}(b) \geq 0$: Negative aFRR capacity bid for block $b$.
    \item $y_{ch}(t) \in \{0,1\}$: Binary variable, 1 if charging in interval $t$, 0 otherwise.
    \item $y_{dis}(t) \in \{0,1\}$: Binary variable, 1 if discharging in interval $t$, 0 otherwise.
\end{itemize}
